\begin{workedexample}[Worked Example: Wavelength and velocity in a substrate]
A signal propagates through bulk PTFE, $\varepsilon_r = 2.1$ (a low-loss
microwave dielectric). The wave slows by $\sqrt{\varepsilon_r}$, so the
\emph{velocity factor} is
\[ \mathrm{VF} = \frac{v}{c} = \frac{1}{\sqrt{\varepsilon_r}} = \frac{1}{\sqrt{2.1}} = 0.69, \]
i.e.\ the wave travels at $69\%$ of the free-space speed. At $5\ \mathrm{GHz}$
the wavelength inside the material is
\[ \lambda = \frac{c}{f\sqrt{\varepsilon_r}} = \frac{3\times10^{8}}{(5\times10^{9})(1.449)} = 41.4\ \mathrm{mm}, \]
versus $60\ \mathrm{mm}$ in free space. (On a real microstrip the relevant
quantity is the \emph{effective} permittivity, treated in Part~III.)
\end{workedexample}

\begin{keytakeaways}
\begin{itemize}[leftmargin=1.4em,itemsep=2pt]
\item A dielectric slows waves: $v = c/\sqrt{\varepsilon_r}$, so wavelength shrinks by $\sqrt{\varepsilon_r}$.
\item The loss tangent $\tan\delta = \varepsilon''/\varepsilon'$ measures dielectric loss;
smaller is better, and it matters more as frequency rises.
\item Both $\varepsilon_r$ and $\tan\delta$ drift with frequency (dispersion), so
substrates are specified at a stated test frequency.
\end{itemize}
\end{keytakeaways}

\begin{exercises}
\item A coaxial cable uses solid PTFE ($\varepsilon_r = 2.1$). What is its velocity factor?
\item Find the wavelength of a $5\ \mathrm{GHz}$ signal inside bulk Rogers RO4350
($\varepsilon_r = 3.66$).
\item At $10\ \mathrm{GHz}$, which loses less energy per wavelength: FR-4
($\tan\delta \approx 0.02$) or PTFE ($\tan\delta \approx 0.001$)?
\end{exercises}

\furtherreading{Pozar~\cite{pozar} (ch.~1); Ludwig and Bogdanov~\cite{ludwig} (ch.~2).}
